% ============================================================================================
\section{I1: Node-Link Diagram (Owner: Sidhanth Prabhu, BT2024027)}
\label{sec:i1}
\textbf{Graph.} The graph is the bipartite subreddit--domain network of Sec.~\ref{sec:derive}: 671 nodes and 2,350 weighted
edges. Communities come from Louvain modularity maximisation~\cite{blondel2008} (NetworkX, weighted, resolution 1, seed 42).
In every drawing the encoding is identical: \emph{size} $\propto\sqrt{\text{strength}}$ (weighted degree = number of top
posts), \emph{colour} = Louvain community, \emph{shape} = node type (square subreddit, circle domain), edge width
$\propto\sqrt{w}$ and edge opacity 0.12. Labels are placed with \texttt{adjustText} to avoid overlaps. Script:
\fpath{infovis/i1_network/run.py}.

\subsection{Network structure}
\figfull{Fig_25_network_statistics.png}{(a) Complementary cumulative degree distribution of the two node types (log--log);
(b) edge-weight histogram; (c) Louvain community sizes split into domains and subreddits (number = subreddits).}{fig:25}
\input{generated/tab_graph}
\input{generated/tab_communities}

\textbf{Inferences.}
(1) The degree distributions of the two node types differ completely (Fig.~\ref{fig:25}a). Domains have mean degree 3.7,
17.8\,\% of them are linked from a single subreddit, and no domain is linked from more than 27 subreddits. Subreddits have
mean degree 54.7 and a heavy tail led by the news-type hubs: UpliftingNews (227), Futurology (182), nottheonion (177),
todayilearned (167), books (161) and worldnews (158). The most shared domains are imgur (27 subreddits), cnn.com and
bbc.com (17 each), and youtube.com and nytimes.com (16). The graph is \emph{hub-and-spoke}: degree
assortativity is $-0.77$, meaning high-degree subreddits link mostly to low-degree domains.
(2) Edge weights are very skewed (median 2, maximum 893 posts, Fig.~\ref{fig:25}b). The five heaviest edges are videos$\rightarrow$youtube (893),
DIY$\rightarrow$imgur (892), listentothis$\rightarrow$youtube (846), Documentaries$\rightarrow$youtube (816) and
gifs$\rightarrow$imgur (545). Together they carry 22\,\% of the total weight of 18,245 posts.
(3) Louvain finds six communities with high modularity ($Q=0.611$), and they are \emph{semantically coherent}
(Table~\ref{tab:comm}). C0 contains the news and science-news subreddits and 391 news outlets (the Guardian, CNN, Business
Insider, Wikipedia, BBC). C1 contains 21 image and GIF communities around imgur, gfycat and giphy. C2 contains the video and
music subreddits around YouTube, Vimeo and SoundCloud. C3 contains movies and television around trade press (Variety,
Hollywood Reporter, Deadline). C4 is r/science alone with 74 journal and science-press domains. C5 groups space, anime and
long-form text communities around twitter.com, NASA and github.io.
(4) The single giant component (diameter 6) holds because a few platforms (imgur, YouTube, Wikipedia, Twitter) are shared
across topics. Without them the graph would split by topic. This shows the structural role of general-purpose hosts as
bridges.

\subsection{Five layouts on the same graph}
We compare five layouts: Fruchterman--Reingold~\cite{fruchterman1991}; Kamada--Kawai~\cite{kamada1989} with edge length
$1/w$; ForceAtlas2~\cite{jacomy2014} (NetworkX 3.7 implementation, scaling 2, gravity 1, distributed attraction, 600
iterations); a spectral embedding from the normalised Laplacian~\cite{shi2000,ng2001}; and a circular layout grouped by
community (Fig.~\ref{fig:26}). The force models are:
\begin{align}
\text{FR:}\ & F_a(d)=d^2/k,\quad F_r(d)=-k^2/d,\quad k=1.6/\sqrt{|V|};\\
\text{KK:}\ & E=\textstyle\sum_{i<j}\tfrac{1}{2}k_{ij}(\lVert p_i-p_j\rVert-\ell\,d_{ij})^2,\ k_{ij}=d_{ij}^{-2};\\
\text{FA2:}\ & F_a=w_{ij}\,d/\deg(i),\quad F_r=k_r\frac{(\deg i+1)(\deg j+1)}{d},
\end{align}
where $d_{ij}$ is the graph distance with weights $1/w$, and ``distributed attraction'' divides the attraction by the
degree, pushing hubs outward. The spectral layout uses eigenvectors 2--3 of
$L_{sym}=I-D^{-1/2}WD^{-1/2}$, rescaled by $D^{-1/2}$. Table~\ref{tab:layouts} and Fig.~\ref{fig:27} quantify each layout:
runtime, exact count of crossing edge pairs, coefficient of variation of edge length, the silhouette of the 2-D positions
grouped by community (a measure of visual cluster separation) and near-overlaps.

\figfull{Fig_26_network_layouts.png}{The same 671-node graph with five layouts and identical encoding (size =
$\sqrt{\text{strength}}$, colour = Louvain community, squares = subreddits, circles = domains, labels = eight strongest
nodes). Runtime under each title.}{fig:26}
\figfull{Fig_27_network_layout_metrics.png}{Quantitative layout comparison: (a) runtime (log scale); (b) edge crossings; (c)
coefficient of variation of edge length; (d) silhouette of node positions by Louvain community. Grey = NetworkX's
unnormalised spectral layout, not drawn in Fig.~\ref{fig:26}.}{fig:27}
\input{generated/tab_layouts}

\textbf{Inferences.}
(1) \textbf{Fruchterman--Reingold} (4.7\,s) gives the most readable overview. The news community C0 forms a large blue cloud
whose outlets surround worldnews and technology. Science (C4, purple) sticks out as its own lobe to the right, and the
image, video and entertainment hubs (C1--C3) stack in a central column. It has the most uniform edge lengths (CV 0.51) and a
positive silhouette (0.12).
(2) \textbf{Kamada--Kawai} is the slowest by two orders of magnitude (526\,s, because it needs all-pairs shortest paths and
$O(|V|^2)$ energy terms). It places the 628 domains on concentric rings by graph distance around a tight core of hubs. The
rings are an artefact of the bipartite distance structure (every domain is at distance 1 or 3 from the core). The result
has the most crossings among the force-based layouts (382,328) and a negative silhouette ($-0.15$), so community
structure is lost.
(3) \textbf{ForceAtlas2} (13.3\,s) with distributed attraction pushes the large subreddit hubs to the \emph{periphery}, each
surrounded by its own exclusive domains (r/science at the top with its 74 science outlets, the news subreddits around the
edge). Shared domains (imgur, YouTube) are pulled to the centre. This ``authorities in the middle, hubs outside'' picture is
the most informative for asking which domains are shared. It has few crossings (226,342), but because shared domains mix in
the centre its silhouette is near zero ($-0.02$).
(4) \textbf{Spectral} (normalised) is the fastest real layout (0.08\,s) and separates communities best (silhouette 0.49,
fewest crossings 220,414). It arranges the communities as \emph{arms}: video and music at the top, images and DIY to the
right, news at the bottom left, and science at the bottom tip. The cost is 469 near-overlapping nodes along each arm, so
individual domains cannot be read. NetworkX's \emph{unnormalised} spectral layout collapses: 666 near-overlaps, CV 12.1,
silhouette $-0.11$. Its Fiedler vectors localise on a few weakly attached nodes, which is why the normalised Laplacian is
required on hub-dominated graphs.
(5) \textbf{Circular} grouped by community is instant and orders nodes perfectly (silhouette 0.19 by construction), but it
has 4.3$\times$ more crossings (966,270) than the force layouts. Inter-community links become a dense hairball of chords,
mostly from the news arc to imgur and YouTube.
(6) Overall: FR for a first look, FA2 to study shared versus exclusive sources, spectral for community separation, and
KK not at all at this size. No layout reduces crossings below about $2.2\times10^5$, which is inherent to a graph with
hubs of degree over 200.

\subsection{Node filtering}
\figfull{Fig_28_network_filtering.png}{Filtering with ForceAtlas2. (a) Full graph (671 nodes, 2,350 edges); (b) after removing
edges with weight $<2$ and then keeping the 2-core (269 nodes, 1,003 edges). Colours are the communities of the unfiltered
graph, and 16 labels are shown on the filtered view.}{fig:28}

\textbf{Inferences.}
(1) The filter removes 60\,\% of nodes and 57\,\% of edges: one-off links (weight 1, 42.4\,\% of edges) and domains linked from a
single community. The remaining backbone is readable enough for 16 labels without overlap.
(2) Removing the spoke domains reveals the macro-structure (Fig.~\ref{fig:28}b). The news hubs (worldnews, news,
nottheonion, technology, Futurology, UpliftingNews, todayilearned) share a dense core of outlets. A chain of image
subreddits (gifs, DIY and others) hangs off imgur and gfycat. Video and music (videos, Documentaries, listentothis) attach
through youtube.com, and television connects the entertainment trade press.
(3) The filter also hides information. r/science's 74 specialised outlets mostly disappear because few of them are shared,
which makes science look more peripheral than it is in Fig.~\ref{fig:26}. The threshold should always be reported.

\subsection{Subreddit similarity projection}
To compare communities directly we project the bipartite graph onto subreddits. Each subreddit is a vector of
TF-IDF-weighted link counts, $x_{sd}=\log(1+w_{sd})\cdot\log(|S|/n_d)$, which down-weights ubiquitous hosts such as imgur.
An edge is drawn where the cosine similarity is at least 0.15 (Fig.~\ref{fig:29}). The positions come from Kamada--Kawai
on distance $1-\cos$, so closeness on the page means similar link profiles.

\figfull{Fig_29_subreddit_projection.png}{One-mode projection: 43 subreddits, an edge where the TF-IDF cosine similarity of
linked domains is at least 0.15 (266 edges, 5 Louvain communities). Colour = A1 topic category; node size = number of top
posts linking externally.}{fig:29}

\textbf{Inferences.}
(1) Two large blocs emerge. One is an \emph{image-sharing} bloc at the top left: pics, aww, gifs, funny, memes,
OldSchoolCool, interestingasfuck, creepy, Art, DIY, food and GetMotivated. The other is a \emph{news/link} bloc at the
bottom right: news, worldnews, nottheonion, technology, Futurology, UpliftingNews and todayilearned, together with
movies, television, books and Music, which link to the same news and trade outlets.
(2) Topic labels and link behaviour disagree in instructive ways. gaming and sports (Entertainment) sit in the image bloc,
because their top posts are images and clips. todayilearned (Knowledge) sits with the news bloc because it links
Wikipedia and newspapers.
(3) The video subreddits (videos, Documentaries, listentothis) form a bridge between the blocs through YouTube and Vimeo.
science, askscience, InternetIsBeautiful and dataisbeautiful are peripheral, each with a distinctive set of sources.
(4) Some similarities of 1.0 (memes, Jokes, WritingPrompts, mildlyinteresting) come from very sparse vectors: these
communities link out only a few times, all to the same host. The projection should be read together with node size,
which is small for them.

\subsection{Interactive version and Gephi export}
\figfull{Fig_30_network_interactive.png}{The interactive PyVis/vis-network page (\fpath{interactive/network.html}).
Physics is ForceAtlas2-based. Interactions: hover tooltips (type, topic or hosting class, posts, partners, community),
scroll zoom, pan, node drag, ``Select a Node'' to highlight a neighbourhood, and a filter by any node attribute (type,
community, category), plus navigation buttons.}{fig:30}
The page (Fig.~\ref{fig:30})~\cite{pyvis} supports all the interactions the assignment requires. With the filter, for
example, \texttt{type = subreddit} isolates the 43 hubs, and \texttt{community = C4} shows r/science's exclusive outlets.
The same graph and a filtered version are exported as GEXF with all node attributes. The README gives step-by-step
instructions to reproduce the ForceAtlas2, Fruchterman--Reingold and Yifan Hu layouts in Gephi; Appendix~\ref{app:gephi}
holds the slots for those screenshots.

\subsection{Observations and limitations (I1)}
The network shows that Reddit's largest communities fall into link ecosystems (news outlets, image hosts, video
platforms, trade press, science press) that only partly match their topics. A few general-purpose platforms hold the whole
graph together. Limitations: the network is built from \emph{top} posts only, so it describes what succeeds, not what is
posted. The $\ge3$-post domain threshold removes the long tail of rare sources. Louvain is stochastic. Across seeds 0--9 it finds 6--8 communities with $Q=0.611$--$0.616$, and the adjusted
Rand index against seed 0 is 0.77--0.95. The large blocs are stable; the small communities (C5) split or merge.
